\documentclass[12pt]{article}
\usepackage{amsmath, amssymb}
\usepackage[margin=2.5cm]{geometry}
\usepackage{parskip}
\usepackage{xcolor}
\usepackage{booktabs}
\usepackage{array}

\title{Likelihood from Hamiltonian Flow: Derivation}
\author{}
\date{}

\begin{document}
\maketitle

\section*{Step 1: Initial Distribution}

Stars are assumed to be drawn at $t=0$ from a 2D Gaussian centred on the midplane:
\[
  f_0(z_0,\, w_0)
  = \frac{1}{2\pi\,\sigma_z\,\sigma_w}
    \exp\!\left(-\frac{z_0^{2}}{2\sigma_z^{2}} - \frac{w_0^{2}}{2\sigma_w^{2}}\right).
\]

\section*{Step 2: Liouville's Theorem}

Hamiltonian dynamics preserve phase-space volume: the Jacobian of any canonical
flow $\Phi^t$ is exactly $1$.  Therefore the probability density at observation
time $t$ is obtained by \emph{back-integrating} each observed point
$(z_\mathrm{obs},w_\mathrm{obs})$ to its initial condition
$(z_0, w_0) = \Phi^{-t}_{h_{dm}}(z_\mathrm{obs}, w_\mathrm{obs})$:
\[
  f_t\!\left(z_\mathrm{obs},\, w_\mathrm{obs}\right)
  = f_0\!\Bigl(\underbrace{\Phi^{-t}_{h_{dm}}(z_\mathrm{obs},\, w_\mathrm{obs})}_{%
      (z_0,\,w_0)\;\text{back-integrated}}\Bigr).
\]
No Jacobian factor appears — it cancels exactly.

\section*{Step 3: The Likelihood}

For $N$ independently observed stars:
\[
\boxed{
  \mathcal{L}(h_{dm})
  = \prod_{i=1}^{N}
    \frac{1}{2\pi\,\sigma_z\,\sigma_w}
    \exp\!\left(-\frac{\bigl[z_0^{(i)}(h_{dm})\bigr]^{2}}{2\sigma_z^{2}}
                -\frac{\bigl[w_0^{(i)}(h_{dm})\bigr]^{2}}{2\sigma_w^{2}}\right)
}
\]
\textbf{Key point:} $z_0^{(i)}$ and $w_0^{(i)}$ are the \emph{back-integrated initial conditions},
not the observed values.  $h_{dm}$ enters exclusively through this back-integration.

\section*{Step 4: Log-Likelihood and Sensitivity}

Taking the log:
\[
  \ln\mathcal{L}(h_{dm})
  = \mathrm{const}
    - \sum_{i=1}^{N}
      \left[
        \frac{\bigl[z_0^{(i)}(h_{dm})\bigr]^{2}}{2\sigma_z^{2}}
        +
        \frac{\bigl[w_0^{(i)}(h_{dm})\bigr]^{2}}{2\sigma_w^{2}}
      \right].
\]
You maximise by finding $h_{dm}$ such that back-integrated orbits are most
consistent with starting near the midplane (small $z_0$, small $w_0$).

The gradient with respect to $h_{dm}$ is:
\[
  \frac{\partial \ln\mathcal{L}}{\partial h_{dm}}
  = -\sum_{i=1}^{N}
    \left[
      \frac{z_0^{(i)}}{\sigma_z^{2}}\,\frac{\partial z_0^{(i)}}{\partial h_{dm}}
      +
      \frac{w_0^{(i)}}{\sigma_w^{2}}\,\frac{\partial w_0^{(i)}}{\partial h_{dm}}
    \right].
\]
The partial derivatives $\partial z_0/\partial h_{dm}$ and $\partial w_0/\partial h_{dm}$
are small because dark matter contributes only $\sim 8\%$ of the total gravitational
force at typical orbit heights --- the back-integrated orbit barely changes when
$h_{dm}$ is varied.

\section*{Step 5: Force Budget}

The gravitational force is a sum over sech$^2$ components:
\[
  K_z(z) = -\frac{d\Phi}{dz}
  = -4\pi G\sum_\alpha \rho_{0,\alpha}\,h_\alpha\,\tanh\!\left(\frac{z}{h_\alpha}\right),
\]
with parameters:

\vspace{0.5em}
\begin{center}
\begin{tabular}{lcccc}
\toprule
Component & $\rho_0\;[M_\odot/\mathrm{pc}^3]$ & $h\;[\mathrm{pc}]$
          & $\rho_0 h$ & DM fraction at $z=250\;\mathrm{pc}$ \\
\midrule
Thin disk  & 0.04 & 200 & 8.0  & \\
Thick disk & 0.01 & 400 & 4.0  & \\
Gas        & 0.05 & 800 & 40.0 & \\
\textbf{DM} & \textbf{0.01} & \textbf{250} & \textbf{2.5} & $\mathbf{8.3\%}$ \\
\midrule
Total      &      &     & 54.5 & \\
\bottomrule
\end{tabular}
\end{center}
\vspace{0.5em}

Gas alone ($\rho_0 h = 40$) is $16\times$ larger than dark matter ($\rho_0 h = 2.5$).
DM never exceeds $\sim 9\%$ of the total force at any relevant height.

\section*{Step 6: Why the $\sigma$ Mismatch Fails}

With a wrong $\sigma_{z,\mathrm{wrong}} < \sigma_{z,\mathrm{true}}$, the biased log-likelihood is:
\[
  \ln\mathcal{L}_\mathrm{biased}(h_{dm})
  = \mathrm{const}
    - \sum_{i}
      \left[
        \frac{\bigl[z_0^{(i)}\bigr]^{2}}{2\,\sigma_{z,\mathrm{wrong}}^{2}}
        +
        \frac{\bigl[w_0^{(i)}\bigr]^{2}}{2\,\sigma_{w,\mathrm{true}}^{2}}
      \right].
\]
Shrinking $\sigma_z$ \emph{over-weights} the $z_0^2$ term, tilting the landscape.
The shift of the peak scales as:
\[
  \Delta h_{dm} \;\propto\;
  \frac{\text{tilt from mismatch}}{\text{curvature of }\ln\mathcal{L}}
  \;\propto\;
  \frac{1/\sigma_{z,\mathrm{wrong}}^{2} - 1/\sigma_{z,\mathrm{true}}^{2}}
       {\;\partial^2 \ln\mathcal{L}/\partial h_{dm}^2\;}.
\]
Since the curvature (Fisher information in $h_{dm}$) is tiny due to the $\sim 8\%$
DM fraction, even a large mismatch in $\sigma_z$ barely moves the optimum.

\end{document}
